\subsection{Construção da rede dirigida e ponderada}
\label{sec:construcao-rede}

A rede tem 67 nós, um por atividade da MIP, identificados pelos códigos do IBGE. Cada nó representa uma origem emissora nas ligações de saída e o setor de um produto final nas ligações de entrada, conforme a representação de \textcite{huang2019}. Em cada nó de destino, a atribuição reúne as emissões associadas aos diferentes usos finais dos produtos da atividade, como consumo das famílias e exportações. A matriz de adjacência ponderada desta rede é definida como:
\begin{equation}
W_{ij}=\begin{cases}
F_{ij}, & i\ne j,\\
0, & i=j.
\end{cases}
\label{eq:rede-sem-diagonal}
\end{equation}
Cada $W_{ij}>0$ define uma aresta $i\rightarrow j$, cujo peso, em Gg de CO\textsubscript{2}, indica a emissão em $i$ atribuída ao produto final de $j$. As duas direções de um par são distintas. Como os pesos incorporam requisitos diretos e indiretos, concatenar arestas não representa uma sequência de etapas produtivas.

A diagonal $F_{ii}$ reúne emissões de $i$ atribuídas à demanda final de seus próprios produtos, incluindo requisitos indiretos. Ela permanece em $F$ e nas contas $c$ e $u$, mas é retirada de $W$, cujo peso total corresponde apenas às atribuições entre atividades distintas.

\subsection{Forças e concentração das relações}
\label{sec:volume-concentracao}

Para examinar como as emissões se distribuem entre atividades emissoras e setores de produtos finais, analisam-se as forças de saída e de entrada\footnote{A força é a soma dos pesos das ligações de um nó em uma dada direção. Neste trabalho, corresponde ao volume intersetorial de emissões, em Gg de CO\textsubscript{2}.} e a concentração\footnote{Indica em que medida o volume atribuído está associado a poucas atividades. Quanto mais próximo de 1, maior a concentração.} das relações representadas em $W$. A análise conjunta permite distinguir atividades com forças semelhantes, mas com diferentes distribuições dos pesos entre destinos, na perspectiva de saída, ou entre origens, na perspectiva de entrada \parencite{huang2019,serrano2009}.

\paragraph{Forças.}
As forças de saída e de entrada de cada atividade são calculadas pela soma dos pesos das respectivas ligações:
\begin{equation}
\begin{aligned}
s_i^{\mathrm{out}}&=\sum_{j\ne i}W_{ij}=c_i-F_{ii},\\
s_i^{\mathrm{in}}&=\sum_{j\ne i}W_{ji}=u_i-F_{ii}.
\end{aligned}
\label{eq:forcas-externas}
\end{equation}
A força de saída $s_i^{\mathrm{out}}$ mede as emissões próprias de $i$ atribuídas aos produtos finais de outras atividades; a força de entrada $s_i^{\mathrm{in}}$ mede as emissões de outras atividades atribuídas aos produtos finais de $i$. A diferença entre força de entrada e força de saída reproduz o saldo entre as contas: $s_i^{\mathrm{in}}-s_i^{\mathrm{out}}=u_i-c_i$.

\paragraph{Concentração.}
Para avaliar a concentração, calculam-se primeiro as participações de cada ligação na força da atividade, separadamente para saída e entrada. Para forças positivas, essas participações são:
\begin{equation}
p_{ij}^{\mathrm{out}}=\frac{W_{ij}}{s_i^{\mathrm{out}}},\qquad
p_{ji}^{\mathrm{in}}=\frac{W_{ji}}{s_i^{\mathrm{in}}},\quad j\ne i.
\label{eq:participacoes-externas}
\end{equation}
As participações somam um em cada perfil de saída ou de entrada. A concentração de cada perfil é medida pela disparidade, denotada por $Y$, que corresponde à soma dos quadrados dessas participações:
\begin{equation}
Y_i^{\mathrm{out}}=\sum_{j\ne i}(p_{ij}^{\mathrm{out}})^2,\qquad
Y_i^{\mathrm{in}}=\sum_{j\ne i}(p_{ji}^{\mathrm{in}})^2.
\label{eq:disparidade}
\end{equation}
Quanto maior $Y$, maior a concentração. Para $k$ relações de igual peso, $Y=1/k$; quando todo o peso está associado a uma única relação, $Y=1$.

Para interpretar essa concentração em termos de quantidade de atividades, calcula-se o número efetivo de contrapartes, denotado por $D_2$, como o inverso da disparidade:
\[
D_{2,i}^{\mathrm{out}}=1/Y_i^{\mathrm{out}},\qquad
D_{2,i}^{\mathrm{in}}=1/Y_i^{\mathrm{in}}.
\]
O número efetivo corresponde à quantidade de contrapartes com participações iguais que produziria a mesma concentração observada. Para $k$ pesos iguais, $D_2=k$; quando uma relação domina, $D_2$ aproxima-se de um. Portanto, valores menores de $D_2$ indicam maior concentração. As participações, a disparidade e o número efetivo são adimensionais e ficam indefinidos quando a força é zero, sem excluir o nó.

Os pesos e as participações identificam as contrapartes, cuja identidade não é informada por $D_2$. O grau conta as arestas positivas de entrada ou saída; a densidade é o número de arestas dividido por $67\times66$. Esses indicadores descrevem a presença de relações, enquanto a disparidade considera a distribuição dos pesos.

\subsection{Extração e avaliação do backbone}
\label{sec:backbone}

O filtro de disparidade é um método que reduz a rede ao selecionar ligações com participações excepcionalmente altas nos pesos de uma origem ou de um destino, segundo um modelo nulo \parencite{serrano2009,huang2019}. A sub-rede formada pelas ligações selecionadas é denominada backbone e preserva seus pesos originais, sem alterar as contas ambientais completas.

\paragraph{Critério de seleção.}
O modelo nulo distribui a força entre $k$ ligações dividindo o intervalo unitário por $k-1$ pontos uniformemente escolhidos. Para $k>1$, a probabilidade de uma parcela atingir ou superar $p$ é $(1-p)^{k-1}$. Para a aresta $i\rightarrow j$, as caudas nas perspectivas de saída e entrada são
\begin{equation}
\begin{aligned}
a_{ij}^{\mathrm{out}}&=\left(1-\frac{W_{ij}}{s_i^{\mathrm{out}}}\right)^{k_i^{\mathrm{out}}-1},\\
a_{ij}^{\mathrm{in}}&=\left(1-\frac{W_{ij}}{s_j^{\mathrm{in}}}\right)^{k_j^{\mathrm{in}}-1}.
\end{aligned}
\label{eq:caudas-filtro}
\end{equation}
Os graus $k$ e as forças são os de $W$. A aresta é selecionada se $a_{ij}^{\mathrm{out}}<\alpha$ \textbf{ou} $a_{ij}^{\mathrm{in}}<\alpha$, registrando-se a extremidade responsável pela seleção. A saída destaca um destino nas atribuições da origem; a entrada destaca uma origem na composição do destino.\footnote{As caudas são calculadas apenas para graus pertinentes maiores que um. Se uma extremidade tem grau um, vale o teste da outra; se ambas têm grau um, a aresta é preservada por convenção de manutenção da relação exclusiva, inclusive em díades isoladas. Essa extensão dos casos degenerados do suplemento de Serrano et al. é registrada separadamente da seleção estatística.}

Adota-se $\alpha=0{,}05$, com $0{,}01$ e $0{,}10$ para sensibilidade. O filtro é aplicado uma vez a $W$ para cada valor, sem atualização dos perfis. A matriz $B_\alpha$ preserva os pesos originais nas arestas selecionadas e zero nas demais, sem renormalização ou corte absoluto adicional.

\paragraph{Avaliação da seleção.}
Comparam-se a fração de arestas retidas e a cobertura dos pesos, globalmente e por atividade:
\begin{equation}
\begin{aligned}
r_\alpha&=\frac{\sum_{ij}(B_\alpha)_{ij}}{\sum_{ij}W_{ij}},\\
r_{\alpha,i}^{\mathrm{out}}&=\frac{\sum_j(B_\alpha)_{ij}}{s_i^{\mathrm{out}}},\qquad
r_{\alpha,i}^{\mathrm{in}}=\frac{\sum_j(B_\alpha)_{ji}}{s_i^{\mathrm{in}}}.
\end{aligned}
\label{eq:cobertura-backbone}
\end{equation}
As razões são definidas apenas para denominadores positivos. A cobertura refere-se à parcela intersetorial de $W$ e permite distinguir retenção agregada de perdas por atividade e direção. A conectividade é descrita pelos graus, isolados e componentes fracamente conectados, obtidos ao desconsiderar a direção das arestas.

A sensibilidade é examinada pelas mudanças nas arestas, coberturas e posições setoriais entre os três valores de $\alpha$. O aninhamento dos conjuntos decorre da regra e não constitui evidência independente de robustez. As caudas do modelo nulo não medem incerteza ambiental, e a seleção não estabelece prioridade de mitigação.

Para interpretar os perfis, as contrapartes são ordenadas por peso em $W$ completa. Destacam-se as cinco maiores e a participação conjunta das demais, sem modificar a rede. A leitura dos casos combina essas participações, as contas completas e os motivos de retenção; a sensibilidade verifica também a permanência das principais relações nos três parâmetros.
